\documentclass[11pt,a4paper]{article}
\usepackage{fontspec}
\usepackage{kotex}
\usepackage{geometry}
\usepackage{graphicx}
\usepackage{xcolor}
\usepackage{tcolorbox}
\usepackage{booktabs}
\usepackage{array}
\usepackage{longtable}
\usepackage{colortbl}
\usepackage{fancyhdr}
\usepackage{titlesec}
\usepackage{hyperref}
\usepackage{emoji}
\usepackage{amsmath}
\usepackage{amssymb}
\usepackage{enumitem}

\geometry{left=25mm, right=25mm, top=30mm, bottom=30mm}

\setmainfont{Noto Sans CJK KR}
\setsansfont{Noto Sans CJK KR}
\setmonofont{Noto Sans CJK KR}
\setemojifont{Noto Color Emoji}

\definecolor{primary}{HTML}{1976D2}
\definecolor{darkbrown}{HTML}{3E2723}
\definecolor{mediumbrown}{HTML}{5D4037}
\definecolor{lightbrown}{HTML}{8D6E63}
\definecolor{cream}{HTML}{FFF3E0}
\definecolor{lightgray}{HTML}{F5F5F5}
\definecolor{tableheader}{HTML}{E8E8E8}

\titleformat{\section}
  {\Large\bfseries\color{primary}}
  {\thesection.}{0.5em}{}
\titleformat{\subsection}
  {\large\bfseries\color{darkbrown}}
  {\thesubsection}{0.5em}{}
\titleformat{\subsubsection}
  {\normalsize\bfseries\color{mediumbrown}}
  {\thesubsubsection}{0.5em}{}

\renewcommand{\contentsname}{목차}
\renewcommand{\figurename}{그림}
\renewcommand{\tablename}{표}

\tcbset{
  tipbox/.style={colback=cream,colframe=primary,boxrule=1.5pt,arc=3mm,
    left=5mm,right=5mm,top=3mm,bottom=3mm,fonttitle=\bfseries\color{primary}},
  formulabox/.style={colback=lightgray,colframe=lightgray,boxrule=0pt,arc=2mm,
    left=5mm,right=5mm,top=3mm,bottom=3mm},
  summarybox/.style={colback=white,colframe=primary,boxrule=2pt,arc=4mm,
    left=5mm,right=5mm,top=5mm,bottom=5mm}
}

\hypersetup{colorlinks=true,linkcolor=primary,urlcolor=primary,bookmarks=true,unicode=true}


\pagestyle{fancy}
\fancyhf{}
\fancyhead[R]{\small\color{lightbrown}마이크로비트 Grove I2C 센서 시뮬레이터 활용가이드 v1.0.0.0}
\fancyfoot[C]{\thepage}
\renewcommand{\headrulewidth}{0.4pt}

\begin{document}

\begin{titlepage}
\centering
\vspace*{3cm}
{\Huge\bfseries\color{primary}\emoji{electric-plug} 마이크로비트 Grove I2C 센서 시뮬레이터}\\[0.5cm]
{\Huge\bfseries\color{primary} 활용가이드}
\\[1.5cm]
{\Large 시뮬레이터 버전: v9.2.0.0}
\\[2cm]
\begin{tcolorbox}[summarybox]
\centering
{\large 이 활용가이드는 마이크로비트 Grove I2C 센서 시뮬레이터의 수업 적용 방법을 안내합니다.}
\end{tcolorbox}
\vfill
{\large 사물인터넷과 빅데이터}\\[0.3cm]
{\large 청주교육대학교 컴퓨터교육과}
\vspace{1cm}
\end{titlepage}

\section*{1. 수업 개요}

\begin{center}
\renewcommand{\arraystretch}{1.4}
\begin{tabular}{|p{2.8cm}|p{12.2cm}|}
\hline
\rowcolor{tableheader}
\textbf{항목} & \textbf{내용} \\
\hline
강좌 & 사물인터넷과 빅데이터 \\
\hline
적용 주차 & 2\textasciitilde{}3주차 \\
\hline
자료 역할 & 참고자료 \\
\hline
대상 & 4학년 (프로그래밍 비전공, 초등 예비교사) \\
\hline
수업 시간 & 45분 (1차시) \\
\hline
\end{tabular}
\end{center}

\subsection*{1.1 학습 목표}

$\bullet$ 마이크로비트 내장 센서와 Grove 외장 센서의 차이를 설명할 수 있다

$\bullet$ 프로젝트 목적에 맞는 센서 조합을 탐색하고 선택할 수 있다

$\bullet$ 센서가 측정하는 데이터 항목과 프로젝트 적합성을 평가할 수 있다



\section*{2. 수업 시나리오 (45분)}

\begin{center}
\renewcommand{\arraystretch}{1.4}
\begin{tabular}{|p{1.5cm}|p{1.5cm}|p{8.6cm}|p{3.4cm}|}
\hline
\rowcolor{tableheader}
\textbf{단계} & \textbf{시간} & \textbf{활동} & \textbf{시뮬레이터 활용} \\
\hline
도입 & 5분 & "여러분 스마트폰에는 센서가 몇 개 있을까요?" 질문으로 시작 & — \\
\hline
탐색 & 10분 & 시뮬레이터의 내장/외장 센서 카드를 둘러보며 종류 파악 & 센서 카드 탐색 \\
\hline
실습 & 15분 & "교실 쾌적도 측정" 주제로 필요한 센서 조합 탐색 & 드래그 앤 드롭 연결 \\
\hline
발표 & 10분 & 조별로 선택한 센서 조합과 이유 공유 & 측정 데이터 패널 \\
\hline
정리 & 5분 & 내장 vs 외장 센서 차이점 정리, 다음 주 프로젝트 안내 & — \\
\hline
\end{tabular}
\end{center}


\section*{3. 학습 질문}

\subsection*{초급 (이해)}

$\bullet$ Q: 내장 센서와 외장 센서의 차이는 무엇인가요?

$\bullet$ Q: Grove Shield에 센서를 몇 개까지 연결할 수 있나요?


\subsection*{중급 (적용)}

$\bullet$ Q: '교실 쾌적도 측정' 프로젝트에 어떤 센서 조합이 필요할까요?

$\bullet$ Q: 내장 온도 센서 대신 SHT40을 사용해야 하는 이유는 무엇인가요?


\subsection*{고급 (분석)}

$\bullet$ Q: 센서 수를 늘리면 항상 더 좋은 데이터를 얻을 수 있나요?

$\bullet$ Q: 같은 데이터를 측정하는 두 센서가 있다면 어떤 기준으로 선택하나요?


\subsection*{확장 (창의)}

$\bullet$ Q: 여러분의 일상에서 IoT 센서가 유용할 상황을 3가지 제안해보세요.

$\bullet$ Q: 스마트 교실을 만든다면 어떤 센서 조합이 필요할까요?



\section*{4. 학생 활동지}

\subsection*{활동: 나의 IoT 프로젝트 센서 조합 찾기}

\textbf{목표}: 프로젝트 주제에 맞는 최적 센서 조합을 탐색한다


\begin{center}
\renewcommand{\arraystretch}{1.4}
\begin{tabular}{|p{1.5cm}|p{6.2cm}|p{7.3cm}|}
\hline
\rowcolor{tableheader}
\textbf{단계} & \textbf{활동} & \textbf{기록} \\
\hline
1 & 프로젝트 주제 작성 & 주제: \_\_\_\_\_\_\_\_\_\_\_\_\_\_\_\_\_ \\
\hline
2 & 필요한 데이터 나열 & ①\_\_\_\_\_\_\_ ②\_\_\_\_\_\_\_ ③\_\_\_\_\_\_\_ \\
\hline
3 & 시뮬레이터에서 센서 연결 & 연결한 센서: \_\_\_\_\_\_\_\_\_\_\_\_\_\_\_\_\_ \\
\hline
4 & 측정 데이터 확인 & 측정 항목: \_\_\_\_\_\_\_\_\_\_\_\_\_\_\_\_\_ \\
\hline
5 & 내장/외장 구분 기록 & 내장: \_\_\_\_개, 외장: \_\_\_\_개 \\
\hline
\end{tabular}
\end{center}


\section*{5. 평가 관찰 체크리스트}

\begin{center}
\renewcommand{\arraystretch}{1.4}
\begin{tabular}{|p{3.0cm}|p{4.9cm}|p{3.9cm}|p{3.2cm}|}
\hline
\rowcolor{tableheader}
\textbf{평가 항목} & \textbf{상} & \textbf{중} & \textbf{하} \\
\hline
센서 종류 이해 & 내장/외장 차이 정확히 설명 & 차이를 대략 이해 & 구분 어려움 \\
\hline
센서 조합 선택 & 프로젝트에 최적 조합 선택 & 관련 센서를 일부 선택 & 무작위 선택 \\
\hline
데이터 항목 이해 & 측정 데이터의 의미를 설명 & 항목 이름을 읽음 & 데이터 확인 안 함 \\
\hline
\end{tabular}
\end{center}


\section*{6. 수업 팁}

$\bullet$ \emoji{light-bulb} 처음에는 키워드 추천 버튼을 활용하게 하면 탐색이 수월합니다

$\bullet$ \emoji{light-bulb} 조별로 다른 프로젝트 주제를 배정하면 다양한 센서 조합이 나옵니다

$\bullet$ \emoji{light-bulb} "내장 센서로 충분한가?" 질문을 통해 비용·정밀도 트레이드오프를 논의하세요



\section*{7. 관련 자료}

\begin{center}
\renewcommand{\arraystretch}{1.4}
\begin{tabular}{|p{1.5cm}|p{7.7cm}|p{5.8cm}|}
\hline
\rowcolor{tableheader}
\textbf{\#} & \textbf{자료명} & \textbf{비고} \\
\hline
1 & 마이크로비트Grove센서시뮬레이터\_퀵가이드 & 소프트웨어 사용 중 빠른 참조 \\
\hline
2 & 마이크로비트Grove센서시뮬레이터\_사용설명서 & 전체 기능 상세 \\
\hline
3 & 마이크로비트Grove센서시뮬레이터\_이론해설서 & 배경 이론 \\
\hline
4 & 마이크로비트Grove센서시뮬레이터\_개발참조 & 기술 사양 \\
\hline
\end{tabular}
\end{center}


\vfill
\begin{center}
\small\color{lightbrown}
\textit{마이크로비트 Grove I2C 센서 시뮬레이터 활용가이드}
\end{center}
\end{document}